% Dinâmica: 2 blocos e 1 roldana
% Faro, qui 08 jan 2026 19:56:23   
% Written by Tordar 
% pstricks2pdf.sh 2blocos_1roldana
\documentclass{article}
\pagestyle{empty}

\usepackage{pst-all} 
\usepackage{pstricks-add}
\usepackage{pst-slpe}
\usepackage{graphicx}
\input{apnum}

% Bloco
\newcommand{\bloco}[5][]{% #1=options, #2=cx, #3=cy, #4=totalWidth, #5=height
  % Left frame (half of total width, full height)
  \psframe[linewidth=1pt,linecolor=gray,fillstyle=slope,slopebegin=white,slopeend=gray,slopesteps=15]
    (!#2 #4 2 div sub #3 #5 2 div sub)                     % lower-left of total area
    (!#2 #3 #5 2 div add)                                  % upper-right of left frame (center x, top)  
  % Right frame (half of total width, full height)
  \psframe[linewidth=1pt,linecolor=gray,fillstyle=slope,slopebegin=gray,slopeend=white,slopesteps=15]
    (!#2 #3 #5 2 div sub)                                  % lower-left of right frame (center x, bottom)
    (!#2 #4 2 div add #3 #5 2 div add)                     % upper-right of total area
  \psframe[#1](!#2 #4 2 div sub #3 #5 2 div sub)(!#2 #4 2 div add #3 #5 2 div add) 
}
% Roldana
\newcommand{\roldana}[5]{% x, y, radius, hole-radius, color
 \pscircle[linewidth=0.5mm,linecolor=lightgray,dimen=middle,fillstyle=radslope,slopebegin=gray,slopeend=white,slopesteps=50](#1,#2){#3}
  \ifnum\pdfstrcmp{#4}{0}=0\relax\else
    \pscircle[fillcolor=white,fillstyle=solid,linewidth=0.8pt](#1,#2){#4}
  \fi
  \pscircle[linewidth=1pt,linecolor=#5!50!black](#1,#2){#3}
}

\begin{document}

% Coordenadas do teto: 
\evaldef\xi{2}
\evaldef\xii{8}
\evaldef\xiii{\xi-1}
\evaldef\xiv{\xii+1}
\evaldef\yi{9}
\evaldef\yii{10}
% Roldana: 
\evaldef\ther{1}
\evaldef\xr{5}
\evaldef\yr{7}
\evaldef\xmr{\xr+1.5*\ther}
% Altura e largura do bloco 1: 
\evaldef\wbi{1.5}
\evaldef\hbi{2}
% Altura e largura do bloco 2: 
\evaldef\wbii{1.5}
\evaldef\hbii{2}
% Coordenadas do bloco 1: 
\evaldef\xbi{\xr-\ther}
\evaldef\ybi{\yr-3}
\evaldef\xmi{\xbi-0.8*\wbi}
% Coordenadas do bloco 2: 
\evaldef\xbii{\xr+\ther}
\evaldef\ybii{\yr-5}
\evaldef\xmii{\xbii+0.8*\wbii}

% Let's go:
\begin{pspicture}(0,0)(10,10)
% Teto: 
\psframe[linewidth=0pt,linecolor=lightgray,dimen=inner,fillstyle=hlines*,fillcolor=lightgray,hatchangle=45](\xi,\yi)(\xii,\yii) 
\psline[linewidth=2pt,linecolor=black](\xiii,\yi)(\xiv,\yi) 
% Roldana: 
\roldana{\xr}{\yr}{\ther}{0.05}{gray} 
\psline[linewidth=2pt,linecolor=black](\xr,\yi)(\xr,\yr) 
% Blocos: 
\psline[linewidth=2pt,linecolor=black](\xbi,\yr)(\xbi,\ybi)
\psline[linewidth=2pt,linecolor=black](\xbii,\yr)(\xbii,\ybii)
\bloco[linewidth=2pt,linecolor=black]{\xbi}{\ybi}{\wbi}{\hbi}
\bloco[linewidth=2pt,linecolor=black]{\xbii}{\ybii}{\wbii}{\hbii}
% Texto: 
\rput{0}(\xmi,\ybi){\scalebox{1.5}{$m_1$}}
\rput{0}(\xmii,\ybii){\scalebox{1.5}{$m_2$}}
\rput{0}(\xmr,\yr){\scalebox{1.5}{$m_R$}}
\end{pspicture}

\end{document}
